\chapter{Figuras y algoritmos de detalle}
\label{anexo:figuras}

Este anexo da la notación del modelo, los dos algoritmos que el cuerpo resume, el ciclo de
estimación del motor fiel (vista F4) y el bootstrap paramétrico del estudio chileno, y las dos
vistas de detalle. En el algoritmo \ref{alg:ciclo}, tres reglas son decisiones de diseño: el
grado efectivo se fija por legislador (D-06); el tiempo se normaliza sobre los periodos
servidos (D-10); y las constantes de las búsquedas se toman de la referencia y se declaran en
la configuración de cada módulo. Dos desviaciones conocidas se declaran (D-19): las búsquedas
escalares restauran al salir el mejor estado observado, donde el bucle original puede terminar
en un paso no evaluado, y una restricción de separación que en la referencia es código
inalcanzable no se porta.

La figura \ref{fig:A1} muestra quién posee el estado del motor fiel. CSVLoader crea la
entrada; DWNominate copia de ella en el constructor lo que necesita (votos, coordenadas,
puntos medios y dispersiones) y no retiene el objeto de entrada; posee por valor la
configuración, los metadatos de periodos y la presencia de legisladores, usa las búsquedas
como funciones libres (D-23) y crea el resultado, con el estado terminal, al terminar (D-22).
La CLI construye la configuración y deja apagados sus tres modos de validación, que fijan un
bloque de parámetros; solo el motor moderno los expone como banderas. El interruptor de la
salvaguarda \texttt{GRID} de 2004 está encendido por omisión y \archivo{--wmay-replication} lo
apaga (D-24).

La figura \ref{fig:A2} sigue una estimación con el motor fiel desde el portal en R. La CLI
carga y valida las entradas (inicialización: semillas, $\beta$ y $w_2$), crea el motor con la
configuración y los datos, llama una vez a \texttt{run()} y exporta parámetros terminales
y el resumen fiel mixto (M-02, pendiente de I2).
La nota apunta documentalmente a F4; no es un uso de interacción UML ni añade otro bucle. Solo la carga y la ejecución están dentro de
bloques de captura: si lanzan una excepción, la CLI termina con código 1 y sin salidas; la
construcción del motor y la exportación quedan fuera, y si un archivo de salida no se puede
abrir, la CLI lo omite y termina con código 0, de modo que es la comprobación de archivos del
portal la que rechaza la corrida. El portal acepta el
resultado solo si el código de salida es cero y las coordenadas, los parámetros y el resumen
se pueden leer (D-29), sin comprobar toda su cobertura o esquema. Sella la corrida (D-30);
con el modo wmay, que es el de omisión, añade
\archivo{--wmay-replication} (D-24). El paralelismo con OpenMP es una capacidad del motor; la
corrida comparativa usa un hilo (D-28).

\section*{Notación}
La notación es la del capítulo 2 del documento de tesis. El legislador $i$
ocupa el punto $\mathbf{x}_i$ en $s$ dimensiones; la votación $j$ tiene dos resultados,
$\mathbf{o}_j^{\mathrm{S}}$ y $\mathbf{o}_j^{\mathrm{N}}$, dados por su punto medio y su
separación, con utilidades y probabilidad de votar a favor
\begin{equation}
u_{ij}^{c} = \exp\Bigl(-\sum_{k=1}^{s} w_k^{2} \bigl(x_{ik} - o_{jk}^{c}\bigr)^{2}\Bigr),
\quad c \in \{\mathrm{S}, \mathrm{N}\},
\label{eq:util}
\end{equation}
\begin{equation}
\Pr\bigl(\mathrm{voto}_{ij} = \mathrm{S}\bigr) = \Phi\bigl(\beta \,[u_{ij}^{\mathrm{S}} - u_{ij}^{\mathrm{N}}]\bigr),
\label{eq:prob}
\end{equation}
con $w_1 = 1$, $w_2$ el peso de la segunda dimensión, $\beta$ el parámetro de señal y $\Phi$
la normal estándar. La log-verosimilitud $\ell$ suma el logaritmo de la probabilidad de cada
voto de las votaciones válidas, las de margen de la minoría ${\ge}\,2{,}5$\,\%, y $\ell/N$ es la
verosimilitud por voto. Dos
configuraciones se comparan centrándolas y alineando una con la otra por la transformación
ortogonal óptima, sin cambio de escala; se informan las correlaciones por dimensión $r_1$ y
$r_2$, la distancia media y la razón de amplitud
\begin{equation}
A = \mathrm{rms}(X)/\mathrm{rms}(X_{\mathrm{ref}}), \qquad
\mathrm{rms}(X) = \Bigl({\textstyle\frac{1}{n}} \sum_{i=1}^{n} \Vert \mathbf{x}_i - \bar{\mathbf{x}} \Vert^{2}\Bigr)^{1/2}.
\label{eq:amplitud}
\end{equation}


\begin{algorithm}[htbp]
\caption{Ciclo de estimación del motor fiel.}
\label{alg:ciclo}
\small
\begin{algorithmic}[1]
\Require matrices de votaciones $V$ de $I$ legisladores y $J$ votaciones, valores iniciales $X^{(0)}$, dimensiones $s$, grado $m$, ciclos $T$
\State $w \gets (1, w_2^{(0)}, \beta^{(0)})$; marcar válida cada votación con margen de la minoría ${\ge}\,2{,}5$\,\%
\For{$c = 1, \dots, T$}
  \If{$s \ge 2$} $w_2 \gets$ búsqueda local por pasos desde un paso de $0{,}01$ \EndIf
  \State $\beta \gets$ búsqueda local por pasos desde un paso de $0{,}1$
  \ForAll{votación válida $j$, en paralelo}
     \State proyectar el punto medio $\mathbf{z}_j$ a la bola unitaria
     \State calcular el plano de corte; si $c = 1$ y no hay parámetros precargados, tomar de él $(\mathbf{z}_j, \mathbf{d}_j)$ iniciales y con ellos la polaridad; en los ciclos siguientes, solo registrarlo
     \State $(\mathbf{z}_j, \mathbf{d}_j) \gets$ cinco rondas de búsqueda lineal alternada en separación y punto medio
  \EndFor
  \State evaluar la verosimilitud global
  \ForAll{legislador $i$ con al menos un periodo servido, en paralelo}
     \State grado efectivo según periodos servidos: constante con menos de cinco, lineal con cinco, cuadrático con seis, cúbico con siete o más, sin exceder $m$
     \State término constante inicial por mínimos cuadrados sobre las coordenadas vigentes (en el primer ciclo, los valores iniciales; con menos de cinco periodos, la del primero servido), reescalado a radio 0,75 si queda fuera de la bola
     \State optimizar el término constante con el legislador inmóvil y la restricción de la bola; luego cada término temporal, desde cero y con los anteriores fijos
     \State reconstruir las coordenadas por periodo sobre su tiempo normalizado
  \EndFor
  \State evaluar la verosimilitud global
\EndFor
\State \Return coordenadas, parámetros de votación, $w$, verosimilitud y estadísticos de clasificación
\end{algorithmic}
\end{algorithm}

\begin{algorithm}[htbp]
\caption{Bootstrap paramétrico de medianas de partido.}
\label{alg:bootstrap}
\small
\begin{algorithmic}[1]
\Require panel observado $V$; ajuste observado $(\hat X, \hat Z, \hat D, \hat\beta, \hat w_2)$, con $\hat Z$ y $\hat D$ el punto medio y la separación de cada votación; membresía $M(i,t)$; cortes $t_1 < t_2 < t_3$; réplicas $B$; estado maestro del generador aleatorio
\For{$b = 1, \dots, B$}
  \State derivar el generador aleatorio de la réplica desde el estado maestro
  \State fijar la máscara de redibujo desde el panel observado: celdas no ausentes, votaciones elegibles y separación no nula
  \ForAll{unidad $t$, diputado $i$, votación $j$ en esa máscara}
    \State $p \gets \Phi\!\left(\hat\beta\,(u^{\mathrm{S}}_{ij} - u^{\mathrm{N}}_{ij})\right)$, con las utilidades de la ecuación (\ref{eq:util}) evaluadas en el ajuste observado
    \State $V_b[i,j,t] \gets$ a favor con probabilidad $p$, en contra en otro caso
  \EndFor
  \State valores iniciales de la réplica: estimador estático por unidad sobre $V_b$, con signos armonizados
  \State $X_b \gets$ motor$(V_b)$ desde esos valores iniciales, con el mismo grado, dimensiones, regla y umbral de elegibilidad, ciclos y optimizador; recalcular las votaciones retenidas sobre $V_b$
  \State alinear $X_b$ a $\hat X$ en la primera dimensión por signo y localización
\EndFor
\State fijar las cohortes desde el ajuste observado y reutilizarlas sin cambio en cada réplica
\ForAll{cohorte $c$, partido $p$}
  \State $\delta \gets m_p(t_3) - m_p(t_1)$ sobre $\hat X$; \quad $\Delta_b \gets$ lo mismo sobre $X_b$; \quad $e_b \gets \Delta_b - \delta$
  \State cota superior básica $\gets \delta - q_\alpha(e)$; \quad cota superior percentil $\gets q_{1-\alpha}(\Delta)$
  \State $p_{\mathrm{izq}} \gets (1 + \#\{e_b \le \delta\})/(B+1)$; sesgo $\gets \bar\Delta - \delta$ junto al percentil
  \State elegible si el partido tiene al menos cinco miembros en $t_1$ y en $t_3$
\EndFor
\State ajustar $p_{\mathrm{izq}}$ por Holm dentro de cada familia (cohorte, par de cortes) sobre los partidos elegibles
\State \Return por (cohorte, partido): $\delta$, cotas, $p$ ajustado, indicador de robustez
\end{algorithmic}
\end{algorithm}

La máscara de redibujo observada no es la máscara de ajuste de la réplica. El generador
conserva las celdas fuera de ella; el motor moderno recalcula la proporción minoritaria en
$V_b$ y retiene votaciones según la misma regla y umbral. Una votación elegible observada puede
dejar de serlo en una réplica. Esto corrige la descripción, sin cambiar el bootstrap ni
atribuirle un daño numérico medido. Las semillas por periodo se regeneran y armonizan en cada
réplica; no se reutiliza el contrato de semilla común por legislador de los tres brazos.

\begin{figure}[htbp]
\centering
\makebox[\textwidth][c]{\includegraphics[width=13.5cm]{A1-clases-estado.pdf}}
\caption[A1. Clases y propiedad del estado del motor fiel]{Clases y propiedad del estado del
motor fiel: DWNominate posee la configuración y el estado numérico, usa búsquedas libres (D-23)
y crea el resultado con el estado terminal (D-22).}
\label{fig:A1}
\end{figure}

\begin{figure}[htbp]
\centering
\makebox[\textwidth][c]{\includegraphics[width=13.5cm]{A2-secuencia-estimacion.pdf}}
\caption[A2. Secuencia de una estimación]{Secuencia de una estimación con el motor fiel desde
el portal en R: salida cero y CSV requeridos legibles, sin validar cobertura completa (D-29).}
\label{fig:A2}
\end{figure}

\chapter[Registro completo de decisiones de diseño][Registro de decisiones]{Registro completo de decisiones de diseño}
\label{anexo:decisiones}

D-01 a D-21 provienen del registro del documento de tesis, fechado en su conjunto al 2 de
octubre de 2026; D-22 a D-32 se agregaron en esta etapa al releer el código para las vistas.
La razón y la evidencia de las decisiones que explican las vistas están en el cuadro
\ref{tab:decisiones}, y la sección \ref{sec:registros} las agrupa en nueve registros.

\footnotesize
\begin{longtable}{L{0.8cm}L{3.5cm}L{2.7cm}L{3.2cm}L{1.85cm}}
\caption[Registro de decisiones de diseño.]{Registro de decisiones de diseño, con la
alternativa descartada, la consecuencia y el estado de cada una.}\label{tab:decisiones-completo}\\
\toprule
Id & Decisión & Alternativa descartada & Consecuencia & Estado \\
\midrule
\endfirsthead
\multicolumn{5}{l}{\footnotesize\emph{Cuadro \ref{tab:decisiones-completo}, continuación}}\\
\toprule
Id & Decisión & Alternativa descartada & Consecuencia & Estado \\
\midrule
\endhead
\bottomrule
\endfoot
D-01 & \texttt{DGESDD} de LAPACK de referencia como rutina predeterminada de
descomposición; sin sustitución silenciosa & descomposición de Eigen con sustitución
automática & la configuración se detiene si falta la biblioteca; cada corrida registra la
rutina & vigente \\
D-02 & perfil sin \texttt{fast-math} ni \texttt{march=native} & aritmética relajada &
comparabilidad entre máquinas por verificar, no identidad demostrada &
vigente \\
D-03 & ninguna reducción de punto flotante dentro de regiones paralelas & reducción de OpenMP
sobre la verosimilitud & suma de parciales en orden de índice; identidad de salidas con
varios hilos pendiente de P-01 e I1 & vigente \\
D-04 & ausentes conservados en la geometría del plano de corte, como en la referencia &
filtrarlos de la nube de puntos & el modo anterior queda disponible por variable de entorno y
se registra & vigente \\
D-05 & valores iniciales canónicos por legislador y periodo, generados con el estimador
estático y versionados con el motor & valores iniciales constantes por legislador & los
valores iniciales son parte del instrumento y se regeneran en cada réplica & vigente \\
D-06 & grado efectivo por legislador según periodos servidos; grado lineal como
especificación principal & grado uniforme por panel; cuadrático como principal & coordenadas:
grado efectivo por trayectoria; resumen: histograma por grado
& vigente \\
D-07 & ciclos externos igualados y truncamiento fijo declarado; ninguna afirmación de
convergencia & regla de parada aplicada a ambos motores & prohíbe un requisito de
convergencia; la comparación solo vale a ciclos iguales & vigente \\
D-08 & cada ejecución registra rutina de descomposición, tratamiento de ausencias y marco de
exportación & declarar la configuración solo en la documentación & un resultado sin sello se
regenera en lugar de citarse & vigente \\
D-09 & proyecciones a la bola unitaria conservadas en los tres sitios de la referencia &
eliminarlas & fidelidad procedimental aunque no cambien el acuerdo & vigente \\
D-10 & exportación en el tiempo normalizado de cada legislador, solo periodos servidos &
rejilla global de periodos & el archivo exportado es la configuración ajustada & vigente \\
D-11 & el fiel conserva sus rutinas; búsquedas NLopt solo en el moderno & adoptar NLopt
en el motor fiel & la pregunta de fidelidad y la de mejora se responden con motores
distintos & vigente \\
D-12 & los motores se comparan y nunca se combinan & elegir un mejor motor & toda fila de
resultado nombra motor, evaluador y rutina de descomposición; las variantes híbridas del
repositorio son experimentales y quedan fuera de todo resultado & vigente \\
D-13 & búsqueda escalar local recentrada en el motor moderno & solución condicional exacta
global & evita el colapso de la segunda dimensión en el primer ciclo & vigente \\
D-14 & tres niveles de verificación con un evaluador común & coincidencia de salidas
finales & un defecto que reproduce la primera dimensión no pasa & vigente \\
D-15 & falsación: la decisión de aceptación rechaza cada versión con un defecto real
reintroducido; para RNF01 se propone situar los umbrales entre un valor defectuoso y uno
corregido & umbrales leídos de una corrida que pasa & la prueba puede fallar; el valor del
umbral corresponde al propietario & vigente; umbral propuesto \\
D-16 & ventana igualada de doce unidades como diseño principal; los demás como sensibilidad
& ventana expansiva como principal & un marco; el bootstrap corre sobre él & propuesta,
pendiente de ratificación \\
D-17 & inferencia solo sobre la primera dimensión; estimación en dos & inferencia en ambas &
la segunda dimensión se describe siempre con su advertencia & vigente \\
D-18 & repositorio público con resultados congelados y licencias; generadores del estudio en
el repositorio de investigación & publicar todo el flujo & la reejecución completa en un
solo comando queda pendiente & vigente \\
D-19 & dos desviaciones declaradas de la referencia: las búsquedas escalares restauran al
salir el mejor estado observado; la restricción de separación que en la referencia es
código inalcanzable no se porta & portar el bucle y la restricción literalmente & fidelidad
procedimental con dos excepciones documentadas en el código y en el registro & vigente \\
D-20 & una mejora de rendimiento fiel exige identidad de salidas numéricas a uno y varios
hilos; metadatos variables se validan aparte & aceptar mejoras dentro de
una tolerancia numérica & la medición antes y después compara solo tiempo; un cambio que
altera un bit es un cambio numérico con su propia verificación & propuesta \\
D-21 & línea base de rendimiento reconstruida antes de cualquier cambio, con repeticiones,
el mismo reloj para todos los brazos y la máquina y la compilación registradas & usar los
tiempos publicados como línea base & los tiempos publicados, de una corrida por celda, no
sirven de línea base & propuesta \\
\addlinespace
D-22 & estado numérico concentrado en \texttt{DWNominate}, trazado miembro a miembro a la
referencia; fases como métodos en el orden original & repartir el estado en objetos por fase
& la estructura no es la que un diseño nuevo elegiría; el motor puede detenerse entre fases &
vigente \\
D-23 & búsquedas como funciones libres con su contexto por argumento & métodos del
orquestador & cada búsqueda puede ejecutarse aislada & vigente \\
D-24 & salvaguarda \texttt{GRID} de 2004 tras un interruptor, apagada en el perfil comparado &
portar un solo linaje & el motor fiel se compara con ambos linajes; el sello registra el modo
& vigente \\
D-25 & perfil de replicación para el brazo medido del motor moderno & medirlo con sus opciones
por omisión & fuera del perfil el moderno cambia más que las búsquedas & vigente \\
D-26 & origen para el legislador sin semilla; la rampa anterior, por variable de entorno &
rampa según la posición en el padrón & arranque neutral, como la referencia; la variable no
reproduce bit a bit el binario anterior & vigente \\
D-27 & el portal rechaza sembrar en proceso un panel de varios periodos & sembrar sobre el
primer periodo & exige una semilla por periodo o un archivo explícito & vigente \\
D-28 & el portal fija el número de hilos y exige uno para Fortran y fiel & dejar los hilos al
entorno & toda comparación a un hilo; el sello registra los hilos & vigente \\
D-29 & salida cero, CSV requeridos legibles y directorio inicialmente vacío;
excluidos como valor de retorno & leer lo que haya en el directorio & rechaza archivos
faltantes o previos; cobertura y esquema completos por verificar & vigente \\
D-30 & sello de cada corrida junto al estado & confiar en la documentación & si el motor no
informa su SVD, el sello la infiere de la construcción y la marca como inferida & vigente \\
D-31 & soporte con \texttt{std::sort} y la descomposición espectral de Eigen & port literal de
\texttt{RSORT} y \texttt{RS} & el orden de los empates puede diferir del de \texttt{RSORT} &
vigente; efecto no medido \\
D-32 & el moderno aborta si la LL previa es finita y $\Delta\ell<-10^{-8}$ & continuar y
exportar & inferencia: detectar una violación del contrato de bloque & vigente \\
\end{longtable}
\normalsize

\chapter[Correspondencia entre las rutinas Fortran y los módulos C++][Correspondencia Fortran y C++]{Correspondencia entre las rutinas Fortran y los módulos C++}
\label{anexo:correspondencia}

Este anexo da el catálogo rutina a rutina que la vista F3 resume (cuadro
\ref{tab:correspondencia}) y las piezas que existen solo en una de las dos variantes Fortran
(cuadro \ref{tab:variantes}).

\begin{table}[H]
\centering
\caption{Módulos del motor fiel y la rutina del programa de referencia que cada uno porta.}
\label{tab:correspondencia}
\footnotesize
\begin{tabular}{L{3.3cm}L{5.8cm}L{3.2cm}}
\toprule
Módulo & Responsabilidad & Rutina de referencia \\
\midrule
\archivo{main_cli} & Línea de comandos, detección de periodos, escritura de salidas y
sellos & programa principal (2004); \texttt{dwnom} y el envoltorio standalone (wmay) \\
\archivo{csv_loader} & Lectura del panel y de los valores iniciales; disposición apilada por
legislador y periodo & lectura de datos \\
\texttt{dwnominate} & Orquestador: estado del modelo y ciclo externo con sus cuatro fases &
\texttt{dwnom} \\
\texttt{likelihood} & Verosimilitud global y estadísticos de clasificación, secuencial y
paralela & \texttt{PLOG} \\
\archivo{normal_cdf} & Tabla precalculada de la distribución normal y sus consultas &
\archivo{init_zdf}, tabla ZDF \\
\archivo{grid_optimizer} & Búsqueda escalar local por pasos para $\beta$ y $w_2$ &
\texttt{SIGMAS}, \texttt{WINT} \\
\archivo{rollcall_derivatives} & Verosimilitud y gradientes por votación & \texttt{PROLLC2} \\
\archivo{rollcall_optimizer} & Búsquedas lineales alternadas de separación y punto medio
& \texttt{RCINT2} \\
\archivo{legislator_derivatives} & Verosimilitud, gradientes e información por legislador
sobre la trayectoria & \texttt{PROX} \\
\archivo{optimize_legislators} & Inicialización por mínimos cuadrados y búsqueda término
a término & \texttt{XINT}; \texttt{GRID} (solo 2004; tras un interruptor, D-24) \\
\archivo{cutting_plane} & Clasificación por plano de corte y dirección de mínima varianza
& \texttt{CUTPLANE}, \texttt{SEARCH} \\
\archivo{cutting_point} & Punto de corte óptimo en una dimensión; códigos de voto;
polaridad & \texttt{JAN1PT}, \texttt{JAN11PT} \\
\archivo{simple_ols}, \archivo{sort_utils} & Pseudoinversa por descomposición espectral de
Eigen con el umbral de la referencia; ordenamiento por índices con \texttt{std::sort}, no un
port del algoritmo de \texttt{RSORT} (D-31) & \texttt{REGA} (con \texttt{RS} en 2004,
\texttt{DSYEV} en wmay), \texttt{RSORT} \\
\bottomrule
\end{tabular}
\end{table}

\begin{table}[H]
\centering
\caption{Piezas propias de cada variante Fortran y su tratamiento en el motor fiel.}
\label{tab:variantes}
\footnotesize
\begin{tabular}{L{2.4cm}L{5.0cm}L{4.9cm}}
\toprule
Variante & Piezas propias & En el motor fiel \\
\midrule
Original 2004 & salvaguarda \texttt{GRID}; descomposición \texttt{LSVRR} de IMSL y su
manejo de errores; \texttt{RS} de EISPACK; formato fijo de entrada; estado en bloques
\texttt{COMMON} (\texttt{/XXCOM/}, \texttt{/MINE/}) & \texttt{GRID} portado tras un
interruptor y apagado en el modo wmay (D-24); \texttt{DGESDD} por omisión (D-01); descomposición
espectral de Eigen (D-31) \\
\addlinespace
wmay & \texttt{DGESDD} y \texttt{DSYEV} de LAPACK; acoplamiento para R; envoltorio
standalone; estado en módulos (\archivo{xxcom_mod}, \archivo{mine_mod}, \archivo{zdf_mod}),
con un \texttt{COMMON /AREA1/} residual en \texttt{REGA}; sin \texttt{GRID} & configuración
igualada en el modo wmay: \texttt{DGESDD}, sin \texttt{GRID}; es la referencia de
reproducción y el evaluador común \\
\bottomrule
\end{tabular}
\end{table}

\section*{Dependencias de módulos C++}

Separadas de F3, las inclusiones verificadas de cabeceras se resumen a un nivel de módulo.
Orquestación incluye directamente los cuatro grupos; se listan también las inclusiones
internas, sin confundirlas con correspondencia histórica ni afirmar ocultamiento completo.

\begin{table}[htbp]
\centering\footnotesize
\caption{Adyacencia de inclusiones entre grupos del motor fiel.}
\begin{tabular}{L{3.0cm}L{4.0cm}L{5.4cm}}
\toprule
Cliente & Proveedor & Evidencia en cabeceras \\
\midrule
Orquestación & búsquedas, geometría, modelo, soporte & \archivo{dwnominate.hpp:14--23} \\
Búsquedas & modelo y soporte & \archivo{grid_optimizer.hpp:12--13}; \archivo{rollcall_optimizer.hpp:12}; \archivo{optimize_legislators.hpp:10--12} \\
Modelo & soporte & \archivo{likelihood.hpp:4} \\
Geometría & soporte & \archivo{cutting_plane.hpp:10}; \archivo{cutting_point.hpp:9} \\
\bottomrule
\end{tabular}
\end{table}

\chapter[Construcción, configuración y manejo de errores][Construcción y errores]{Construcción, configuración y manejo de errores}
\label{anexo:construccion}

Este anexo da la vista física del instrumento \parencite[cap.~2]{swebok2025}: cómo se
construyen los tres binarios y el portal en R, qué configuración fija su resultado (cuadro
\ref{tab:construccion}), y cómo responde cada componente a un error \parencite[cap.~3,
§3.5]{swebok2025} (cuadro \ref{tab:errores}). Las rutas siguen la convención del cuadro
\ref{tab:decisiones}.

\begin{table}[H]
\centering
\caption{Construcción y configuración de los tres binarios y del portal en R.}
\label{tab:construccion}
\footnotesize
\begin{tabular}{L{2.3cm}L{3.0cm}L{3.2cm}L{3.6cm}}
\toprule
Componente & Construcción & Dependencias & Configuración que fija el resultado \\
\midrule
Referencia wmay standalone & gfortran; \archivo{DW-NOMINATE-wmay.f} sin cambios, con
\archivo{r_compat.f90} en lugar de las rutinas de impresión y salida de R y
\archivo{standalone_main.f90} & LAPACK y BLAS al enlazar (\texttt{DGESDD}, \texttt{DSYEV}) &
lee y escribe archivos; solo semilla por legislador; modo \texttt{evaluate}; se ejecuta desde
su directorio, que no crea la salida. El original 2004 se compila aparte y es autocontenido \\
\addlinespace
Motor fiel & C++17 con CMake; \texttt{-O3} sin \emph{fast-math} ni \texttt{-march=native} &
OpenMP y Eigen 3.4.0, obligatorios; LAPACK para \texttt{DGESDD}
(\archivo{USE_SYSTEM_DGESDD}, encendida) & sin LAPACK o sin Eigen la construcción se detiene;
la ruta Eigen solo con la opción apagada, y avisa si cae a Jacobi; el resumen registra la SVD
y el marco de exportación; hilos por \archivo{OMP_NUM_THREADS} \\
\addlinespace
Motor moderno & C++20 con CMake; sin \emph{fast-math} & Eigen y NLopt, descargados si no
están instalados; OpenMP opcional, apagado por omisión; LAPACK para \texttt{DGESDD} & misma
detención si falta LAPACK; pruebas con \texttt{BUILD\_TESTING}; \archivo{--threads}, que el
resumen registra \\
\addlinespace
Portal en R (\archivo{dwnom.R}) & R base & \texttt{wnominate} y \texttt{pscl}, solo para
sembrar en proceso & ubica los binarios según la disposición del repositorio; fija hilos,
perfil y semilla (D-24, D-25, D-28); sella la corrida (D-30) \\
\bottomrule
\end{tabular}
\end{table}

\footnotesize
\begin{longtable}{L{4.0cm}L{1.9cm}L{3.6cm}L{2.8cm}}
\caption{Manejo de errores y excepciones, por componente.}\label{tab:errores}\\
\toprule
Condición & Dónde & Respuesta & Evidencia \\
\midrule
\endfirsthead
\multicolumn{4}{l}{\footnotesize\emph{Cuadro \ref{tab:errores}, continuación}}\\
\toprule
Condición & Dónde & Respuesta & Evidencia \\
\midrule
\endhead
\bottomrule
\endfoot
falta LAPACK con \texttt{DGESDD} encendida, o falta Eigen o NLopt & CMake & error de
configuración; no hay binario & F \archivo{CMakeLists.txt:62}, \archivo{:100}; M
\archivo{CMakeLists.txt:44, 58, 76} \\
falta LAPACKE en la ruta de Eigen & CMake, fiel & aviso explícito; se usa Jacobi y el
resumen lo registra & F \archivo{CMakeLists.txt:84} \\
no hay matrices de votos o no se puede crear el directorio de salida & CLI & mensaje y código
1 & F \archivo{main_cli.cpp:460-485} \\
falta o no se lee una matriz o el archivo de semillas & carga & excepción; código 1, sin
salidas & F \archivo{csv_loader.cpp:130, 292}, \archivo{main_cli.cpp:539-547} \\
excepción durante la ejecución, como una falla de \texttt{DGESDD} & motor & código 1, sin
salidas & F \archivo{cutting_plane.cpp:313, 323}, \archivo{main_cli.cpp:583-586} \\
excepción al construir el motor & CLI, fiel & no se captura: la llamada está fuera de todo
bloque \texttt{try} y el proceso termina de forma anormal & F \archivo{main_cli.cpp:576} \\
no se puede abrir un archivo de salida & salida, fiel & lo omite con un mensaje y
termina con código 0; el portal rechaza la corrida por archivo faltante & F
\archivo{main_cli.cpp:228-231}; R 387-391 \\
argumento desconocido & CLI & aviso; la corrida sigue & F \archivo{main_cli.cpp:175} \\
LL previa finita y $\Delta\ell<-10^{-8}$ entre ciclos & motor moderno & excepción; código 1 & M
\archivo{dwnominate.cpp:388-393} \\
binario inexistente; hilos distintos de uno para Fortran o fiel; margen de la minoría distinto de 0,025;
\texttt{periods} distinto del panel; directorio de salida con archivos & portal & detención
antes de invocar el motor & R 84, 263, 268, 274, 279 \\
sembrar en proceso un panel de varios periodos; semilla por fila para wmay; semilla con
identificadores repetidos o coordenadas no finitas & portal & detención antes de invocar el
motor & R 293-299, 304, 320-326 \\
código de salida distinto de cero; faltan coordenadas, parámetros o resumen & portal &
detención; las salidas no se aceptan & R 380-391 \\
legislador que vota y no queda colocado & portal & no es error: fila en \archivo{DROPPED.csv}
y en el valor de retorno & R 394-416 \\
\end{longtable}
\normalsize

\chapter{Contratos de archivos}
\label{anexo:contratos}

Este anexo fija el contrato de archivos de los motores que la vista F6 abstrae (cuadro
\ref{tab:contratos}) y las reglas de resolución de registros que la figura declara con
multiplicidades opcionales.

\begin{table}[H]
\centering
\caption[Contrato de archivos de los motores.]{Contrato de archivos de los motores. Todas
las salidas van al directorio indicado en la línea de comandos; la referencia escribe un
estado equivalente con otros nombres de archivo y de campo.}
\label{tab:contratos}
\footnotesize
\begin{tabular}{L{3.6cm}L{5.2cm}L{3.6cm}}
\toprule
Archivo & Contenido & Observación \\
\midrule
\multicolumn{3}{l}{\emph{Entradas}} \\
\archivo{votes_matrix_p<N>.csv} & fila por legislador, columna por votación; códigos de la
referencia: 1 a 3 a favor, 4 a 6 en contra, 0 o mayor que 6 ausente; celda vacía o ilegible,
ausente & un archivo por periodo, numerados desde 1 sin huecos; el portal escribe 9 como
ausente \\
\archivo{legislator_metadata.csv} & identificador, nombre, partido, región, distrito &
inerte para la estimación \\
\archivo{wnominate_coordinates.csv} & semilla por legislador: identificador y dos coordenadas &
el binario exige el archivo, no cada legislador; único contrato que wmay admite \\
semilla por fila (\archivo{--seed-per-period}) & identificador, periodo, dos coordenadas &
opcional; prioridad sobre la anterior \\
parámetros de votación de referencia & sesión, identificador, punto medio y separación por
dimensión & opcional; si falta, parten en cero \\
\addlinespace
\multicolumn{3}{l}{\emph{Salidas}} \\
\archivo{cpp_coordinates_all_periods.csv} & legislador, periodo, coordenadas, grado
efectivo & solo periodos servidos; marco del optimizador \\
variante \archivo{_corrected} & las mismas coordenadas con los signos invertidos según la
regla de orientación & no es un nuevo ajuste \\
\archivo{cpp_bill_parameters.csv} & votación, punto medio y separación por dimensión &
índice global de votación \\
\archivo{cpp_summary.csv} & verosimilitud, ciclos, votos válidos, clasificación, pesos,
grado pedido, histograma de grado efectivo y tiempo & LL terminal; clasificación de fase
previa (M-02); SVD, ausencias y marco registrados \\
\archivo{RUN-STAMP.csv}, \archivo{DROPPED.csv} & sello de la corrida y excluidos, escritos por
el portal & junto al estado; el portal también los devuelve \\
\bottomrule
\end{tabular}
\end{table}

Reglas de resolución. El padrón y los archivos de semillas son registros crudos: pueden traer
identificadores fuera del panel, y el motor no exige que cada registro calce con exactamente
una fila servida; el portal sí rechaza una semilla con identificadores repetidos. El
emparejamiento de coordenadas entre corridas es por (legislador, periodo), sin validación uno
a uno. La identidad del panel requiere hashes de las matrices, del orden de identificadores
de legisladores, votaciones y periodos, de los archivos de semillas y de la configuración.
Se conserva la procedencia en los manifiestos de linaje y en \archivo{RUN-STAMP.csv}
(binario, semilla, argumentos y backend), sin atribuir al sello hashes que no escribe.
Los recuentos crudos, no ausentes y ajustados son diagnósticos separados; no identifican la matriz.

\chapter{Perfil de ejecución}
\label{anexo:perfil}

Este anexo reúne en el cuadro de perfil de ejecución (cuadro \ref{tab:perfil}) las banderas,
la rutina de SVD, los hilos, el arranque, el evaluador de la verosimilitud, los ciclos y el
contrato de semillas que las vistas F4, T5 y F7 dejan fuera de las cajas, junto con las
invocaciones del protocolo de reproducción wmay. Este cuadro fija esa comparación de tres
brazos; el protocolo original 2004 (cuadro \ref{tab:protocolos}) declara sus diferencias.
El bootstrap regenera semillas por periodo y elegibilidad por réplica (algoritmo \ref{alg:bootstrap}).

\footnotesize
\begin{longtable}{L{2.6cm}L{5.0cm}L{4.8cm}}
\caption[Perfil de ejecución.]{Perfil de ejecución de los motores, con el perfil
comparado.}\label{tab:perfil}\\
\toprule
Aspecto & Opción, variable o invocación & Perfil comparado y observación \\
\midrule
\endfirsthead
\multicolumn{3}{l}{\footnotesize\emph{Cuadro \ref{tab:perfil}, continuación}}\\
\toprule
Aspecto & Opción, variable o invocación & Perfil comparado y observación \\
\midrule
\endhead
\bottomrule
\endfoot
Entradas y salidas & \archivo{--input-dir}, \archivo{--output-dir}, \archivo{--wnominate},
\archivo{--seed-per-period} & código de salida distinto de cero si falta una matriz o el
archivo de semillas (cuadro \ref{tab:errores}) \\
\addlinespace
Estimación & \archivo{--model} (grado), \archivo{--iterations} (ciclos), \archivo{--periods},
\archivo{--dimensions} & grado declarado por panel; cuatro ciclos; dos dimensiones; en una
dimensión no se ajusta el peso y el corte usa \texttt{JAN11PT} \\
\addlinespace
Contrato de semillas & semilla por legislador o por fila & por legislador en todos los
brazos, único contrato que wmay admite; por fila solo como control \\
\addlinespace
Linaje del motor fiel & \archivo{--wmay-replication}; en el portal, el modo wmay o el de la
salvaguarda de 2004 & modo wmay: \texttt{GRID} apagado (D-24) \\
\addlinespace
Perfil del motor moderno & \archivo{--fortran-replication}; el portal añade
\archivo{--optimizer-precision=strict} & fija las opciones canónicas (COBYLA, búsqueda escalar
local, votaciones desde cero, sin identificación adicional ni tolerancias adaptativas, iguales
a las de omisión); cambia la tabla de verosimilitud, la de la referencia en lugar de la normal
continua, y retira las cotas globales de $w_2$ y $\beta$ (D-25) \\
\addlinespace
Búsquedas del motor moderno & \archivo{--scalar-search}, \archivo{--block-solver} & local y
COBYLA en el perfil comparado \\
\addlinespace
Parada opcional del moderno & \archivo{--convergence-abs}, \archivo{--convergence-rel},
\archivo{--convergence-patience}, \archivo{--min-iterations} & desactivada por omisión
(tolerancias en cero); nunca criterio de parada en una comparación (D-07) \\
\addlinespace
Arranque & escalares iniciales de la referencia; votaciones en cero; legislador sin semilla
en el origen & iguales en ambos motores (D-25, D-26) \\
\addlinespace
Evaluador de la verosimilitud & fiel: tabla precalculada de la referencia; moderno:
\archivo{--likelihood} (continua por omisión) & tabla de la referencia en ambos motores (D-25)
\\
\addlinespace
Rutina de SVD & \archivo{USE_SYSTEM_DGESDD} (fiel), \archivo{DWNOMINATE_USE_SYSTEM_DGESDD}
(moderno) & \texttt{DGESDD}; registrada en el resumen del fiel y en el sello (D-01, D-30) \\
\addlinespace
Hilos & \archivo{OMP_NUM_THREADS} (fiel); \archivo{--threads} (moderno) & uno en toda
comparación, impuesto por el portal (D-28) \\
\addlinespace
Modos anteriores del fiel & \archivo{DWNOM_EXPORT_GLOBAL_T},
\archivo{DWNOM_CUTPLANE_FILTER_ABSENT}, \archivo{DWNOM_SEED_FALLBACK_RAMP} & conservan
comportamientos previos a una corrección; los que cambian un número publicado quedan en el
resumen \\
\addlinespace
Volcado de fases & \archivo{DWNOM_PHASE_DUMP} & registro por votación con marcas de fase para
el nivel de componente; inerte en otro caso \\
\addlinespace
Llamada desde el portal & \texttt{dwnom\_run} con directorio de entrada, motor, binarios,
directorio de corridas, panel, grado, ciclos y semillas & devuelve coordenadas, parámetros,
escalares, excluidos y sello; salida cero y CSV requeridos legibles (figura
\ref{fig:A2}) \\
\addlinespace
Reevaluación bajo evaluador común & \archivo{stage_terminal_state.py}; wmay standalone en modo
\texttt{evaluate} & cero ciclos; el arnés exige su directorio de trabajo \\
\addlinespace
Comparación de coordenadas & \archivo{map_agreement.py} & emparejamiento por (legislador,
periodo); alineación ortogonal sin reescalar; marco global o por periodo declarado \\
\end{longtable}
\normalsize

\chapter{Narrativas de los casos de uso}
\label{anexo:narrativas}

Este anexo da, para cada caso de uso de la vista F2, la precondición, la condición de fallo y
el criterio de término que la figura no dibuja (cuadro \ref{tab:narrativas}).

\footnotesize
\begin{longtable}{L{2.6cm}L{3.2cm}L{3.2cm}L{3.2cm}}
\caption{Narrativas de los casos de uso: precondición, fallo y término.}\label{tab:narrativas}\\
\toprule
Caso de uso & Precondición & Fallo & Término \\
\midrule
\endfirsthead
\multicolumn{4}{l}{\footnotesize\emph{Cuadro \ref{tab:narrativas}, continuación}}\\
\toprule
Caso de uso & Precondición & Fallo & Término \\
\midrule
\endhead
\bottomrule
\endfoot
Declarar panel y semillas (RF01, RF02) & directorio de panel con una matriz por periodo y el
archivo de semillas por legislador; opcionalmente semillas por fila y parámetros de votación
de referencia & falta una matriz o el archivo de semillas: error y código de salida distinto
de cero, sin salidas & el motor reporta cuántas coordenadas iniciales aplicó por cada vía,
cuántos legisladores partieron del origen o quedaron excluidos y qué grado efectivo recibió
cada uno; un argumento desconocido y la ausencia del archivo opcional se aceptan con aviso \\
\addlinespace
Estimar estático o dinámico (RF03, RF04) & entradas declaradas; perfil de la corrida
(dimensiones, grado, ciclos) & código de salida distinto de cero, o salida cero con archivos
incompletos: el portal rechaza la corrida & directorio de estado con coordenadas, parámetros
de votación y resumen con su procedencia; código de salida cero \\
\addlinespace
Validar con el caso chileno (RF10 a RF12) & ajuste observado sobre el panel del 55.º
periodo; cortes y cohortes fijados desde el calendario antes de inspeccionar estimación
alguna & una réplica cuyo número de periodos difiera del observado o cuyo grado efectivo
exceda el lineal se rechaza; la orientación de la primera dimensión se verifica antes de
leer réplica alguna & por (cohorte, partido): desplazamiento, cotas, valor $p$ ajustado e
indicador de robustez, con rejilla, grado, valores iniciales y motor declarados \\
\addlinespace
Verificar y comparar (RF08, RF09) & dos corridas controladas que superan el control de
paridad de entradas: mismo panel, semillas, ciclos, grado, códigos y margen de la minoría, marco de
exportación y configuración de compilación & una cláusula de paridad que falla: la
comparación no se ejecuta; una discrepancia con sello distinto se reporta como configuración
distinta, no como defecto & informe de comparación: diferencia de verosimilitud bajo
evaluador común, desacuerdo absoluto por votación, correlaciones por dimensión, razón de
amplitud, peso de la segunda dimensión de ambos motores y reflexiones, con cobertura y marco
declarados \\
\end{longtable}
\normalsize

\chapter[Conjunto de paneles de prueba y ajustes igualados][Paneles de prueba]{Conjunto de paneles de prueba y ajustes igualados}
\label{anexo:paneles}

Este anexo detalla el conjunto de paneles de las pruebas exhaustivas (cuadro
\ref{tab:paneles}) y los estadísticos y ajustes igualados que la sección
\ref{sec:experimental} resume (cuadro \ref{tab:estadisticos}).

\begin{table}[H]
\centering
\caption[Conjunto de paneles de prueba.]{Conjunto de paneles de prueba. La descripción identifica el escenario; los hashes de
matriz, orden, semillas y configuración identifican la entrada. Los recuentos son diagnósticos.}
\label{tab:paneles}
\footnotesize
\begin{tabular}{L{3.4cm}L{3.9cm}L{1.5cm}L{3.5cm}}
\toprule
Panel & Escenario & Grado & Brazos de referencia \\
\midrule
Legislaturas 353, 366 y 368 (estáticos) & una legislatura cada uno; bytes idénticos a un
periodo del panel de 23 periodos de la curación anterior & 0 & referencia, fiel, moderno \\
Senado de EE.\,UU., 90.º Congreso (estático) & un periodo; comparación externa & 0 &
referencia, fiel, moderno \\
Senado de EE.\,UU., Congresos 111 a 115 & cinco periodos; valores iniciales reparados & 1 &
los tres, en dos máquinas \\
Cámara de Diputados 2002--2022 & 23 periodos, extendido a marzo de 2022 & 1 & los tres, en
dos máquinas \\
Cámara de Diputados, curación anterior & 23 periodos, legislaturas 346 a 368 & 2 & los tres \\
55.º periodo legislativo & doce unidades de cuatro meses & 1 & referencia; el brazo fiel
debe congelarse \\
Sintéticos (tres niveles de información en la segunda dimensión) & verdad conocida por
construcción & 0 & archivos de verdad \\
\bottomrule
\end{tabular}
\end{table}

\footnotesize
\begin{longtable}{L{3.5cm}L{8.9cm}}
\caption[Estadísticos y ajustes igualados de las pruebas exhaustivas.]{Estadísticos de las
pruebas exhaustivas, por nivel de verificación, y ajustes del protocolo wmay de tres brazos.
La aceptación contra 2004 registra las diferencias del cuadro \ref{tab:protocolos}; RNF01 queda abierto.}\label{tab:estadisticos}\\
\toprule
Elemento & Especificación \\
\midrule
\endfirsthead
\multicolumn{2}{l}{\footnotesize\emph{Cuadro \ref{tab:estadisticos}, continuación}}\\
\toprule
Elemento & Especificación \\
\midrule
\endhead
\bottomrule
\endfoot
\multicolumn{2}{l}{\emph{Estadísticos, por nivel}} \\
Componente & tasa de error del plano de corte; identidad de los escalares tras el primer
ciclo \\
Estado & diferencia de verosimilitud por voto bajo el evaluador común, con el desacuerdo
absoluto por votación al lado; identidad del recuento de votos válidos \\
Mapa & correlaciones por dimensión en el marco global, razón de amplitud, distancia media y
reflexiones \\
Determinismo & identidad de salidas numéricas con uno, dos, cuatro y doce hilos; metadatos aparte \\
\addlinespace
\multicolumn{2}{l}{\emph{Ajustes igualados}} \\
Ciclos & cuatro, la convención de la comparación publicada, declarados en cada sello \\
Grado & el declarado por panel \\
Valores iniciales & en la comparación wmay, por legislador en todos los brazos, porque wmay solo lee esa
forma; los valores por periodo como control \\
Hilos & uno en toda comparación \\
Motor moderno & modo de replicación y precisión estricta \\
Rutina de descomposición & protocolo wmay: \texttt{DGESDD}; original 2004: \texttt{LSVRR}; fiel: backend registrado \\
Aritmética & estricta, sin \emph{fast-math} \\
Escalares iniciales & los de la referencia \\
Dimensiones & dos \\
\end{longtable}
\normalsize

\chapter{Diseño del estudio chileno}
\label{anexo:chile}

Este anexo da el diseño del estudio chileno que la sección \ref{sec:algoritmos} resume: el
contraste, la inferencia y la periodización con su regla de discretización (tarea E3).

El estudio estima el desplazamiento de las medianas de partido en la primera dimensión durante
el 55.º periodo legislativo, dividido en tres fases por el estallido social del 18 de octubre
de 2019, el plebiscito del 25 de octubre de 2020 y el fin del periodo el 10 de marzo de 2022,
cortes fijados desde el calendario antes de inspeccionar estimación alguna. Para el partido
$p$, el contraste es $\Delta_p = m_{p3} - m_{p1}$, con $m_{pk}$ la mediana de la primera
dimensión de sus diputados activos al fin de la fase $k$, y la hipótesis alternativa
unilateral $\Delta_p < 0$ se fijó antes del análisis \parencite{nieves2026jcc}. La inferencia
sigue el bootstrap paramétrico establecido para este estimador
\parencite{lewis2004measuring,carroll2009measuring}: redibujar cada voto desde su probabilidad
ajustada y reestimar todo el instrumento en cada réplica (algoritmo \ref{alg:bootstrap}). Tres
decisiones lo distinguen del uso habitual: los valores iniciales se regeneran dentro de cada
réplica, porque reutilizar los observados con pocos ciclos puede anclar las soluciones; solo
se redibujan las votaciones válidas; y cada réplica se alinea a la primera dimensión observada
por signo y localización (D-17). Un partido se declara robustamente desplazado solo si ambas
cotas son negativas y el valor $p$ ajustado por \textcite{holm1979simple} es menor que el
nivel (RF11).

La subdivisión del tiempo es una elección del analista que los datos no fijan y que mueve el
resultado (RF12); el diseño la trata como un parámetro declarado (cuadro \ref{tab:disenos}).
El principal, propuesto y pendiente de ratificación (D-16), es la ventana igualada: un solo
ajuste sobre las doce unidades de cuatro meses del periodo, en un solo marco, sobre el cual
corre el bootstrap; los otros tres son análisis de sensibilidad. La regla de discretización
tiene tres cláusulas. El grado es siempre lineal, para que una rejilla más fina no cambie la
mezcla de grados efectivos. La unidad debe ser lo bastante corta para que la cohorte que entró
al inicio alcance cinco unidades servidas, el mínimo del término lineal (D-06), y lo bastante
larga para que ninguna unidad quede vacía en el receso de verano: con el umbral de veinte
votos, la de cuatro meses es la más corta que el registro sostiene. Y toda cifra declara
rejilla, grado y valores iniciales. El modelo se estima en dos dimensiones, pero la
inferencia se hace solo sobre la primera (D-17), porque la segunda es considerablemente más
sensible entre implementaciones, optimizadores, horizontes y réplicas
\parencite{nieves2026jcc}.

\begin{table}[H]
\centering
\caption{Diseños de periodización: qué varía cada uno, cómo trata los marcos y qué pregunta
responde.}
\label{tab:disenos}
\footnotesize
\begin{tabular}{L{2.9cm}L{3.5cm}L{2.8cm}L{3.1cm}}
\toprule
Diseño & Varía o fija & Marcos & Pregunta \\
\midrule
Ventana igualada (principal, pendiente de ratificación) & un ajuste; doce unidades de cuatro
meses; grado lineal & un marco; solo el control de orientación & cuánto se desplazaron las
medianas en un marco común \\
\addlinespace
Ventana expansiva & tres ajustes anidados que terminan en cada frontera & rotación estimada en
las unidades de calentamiento & cuánto depende una posición terminal del registro anterior \\
\addlinespace
Por tramos & tres ajustes independientes, uno por fase, con seis unidades internas cada uno &
tres marcos lado a lado & si hubo un salto en la frontera \\
\addlinespace
Ajuste único rebanado & un ajuste sobre las unidades del diseño por tramos & sin realineación
& control de los dos anteriores \\
\bottomrule
\end{tabular}
\end{table}

\chapter[Mejora de rendimiento: cambios candidatos y protocolo][Mejora de rendimiento]{Mejora de rendimiento: cambios candidatos y protocolo}
\label{anexo:rendimiento}

Este anexo da el catálogo completo de cambios candidatos de la mejora de rendimiento del
motor fiel (cuadro \ref{tab:candidatos}) y el protocolo de medición paso a paso (cuadro
\ref{tab:protocolo}) que la sección \ref{sec:rendimiento} resume. Ningún cambio está
implementado. La inserción concurrente P-01 está en la fase de legisladores: el lazo paralelo
inserta en el contenedor de periodos servidos por legislador sin sección crítica, mientras los
coeficientes temporales tienen una preinicialización en serie hecha precisamente para evitar
ese riesgo (F \archivo{dwnominate.cpp:1463-1465}, \archivo{:1478}, \archivo{:1814}); la corrección
es extender esa preinicialización. Por inspección, la carrera no toca los números ajustados,
pero alimenta el marco de exportación.

\footnotesize
\begin{longtable}{L{0.8cm}L{5.3cm}L{2.2cm}L{4.0cm}}
\caption[Cambios candidatos de rendimiento.]{Cambios candidatos de la mejora de rendimiento
del motor fiel, la fase que afectan y su efecto esperado sobre los resultados. Ninguno está
implementado; el efecto esperado es una hipótesis que la condición de aceptación somete a
prueba.}\label{tab:candidatos}\\
\toprule
Id & Cambio & Fase & Efecto esperado en los resultados \\
\midrule
\endfirsthead
\multicolumn{4}{l}{\footnotesize\emph{Cuadro \ref{tab:candidatos}, continuación}}\\
\toprule
Id & Cambio & Fase & Efecto esperado en los resultados \\
\midrule
\endhead
\bottomrule
\endfoot
M-01 & restringir los lazos de cada votación y las evaluaciones globales al bloque de filas
de su periodo & votaciones; verosimilitud global & ninguno: mismas celdas no ausentes, en el
mismo orden \\
M-02 & guardar los votos por bloques de periodo con índice local, contiguos por votación, con
una vista por legislador para los lazos que la usan & todas; carga & ninguno si cada lazo
conserva su orden; reordenar lazos está excluido \\
M-03 & listas precalculadas, en orden ascendente, de los votantes de cada votación y de las
votaciones de cada legislador & votaciones; verosimilitud global & ninguno: mismo orden de
suma \\
M-04 & eliminar asignaciones y copias dentro de los lazos internos; memoria reutilizable por
hilo & todas & ninguno: mismos valores \\
M-05 & evaluar solo la verosimilitud en los puntos de búsqueda lineal y en la primera y
última evaluación de cada votación & votaciones; legisladores & ninguno esperado; se verifica
porque la primera evaluación inicia la historia de la búsqueda \\
M-06 & una sola consulta a la tabla normal por voto para el valor y su logaritmo &
votaciones; verosimilitud & ninguno: mismas entradas de la tabla \\
M-07 & preparar una sola vez los datos invariantes de cada votación y suprimir el trabajo
muerto del plano de corte & votaciones & ninguno: mismos arreglos, en el mismo orden \\
M-08 & guardar invariantes dentro de las búsquedas escalares & escalares & ninguno si se
respeta el orden de evaluación; valor esperado bajo \\
M-09 & carga en tiempo proporcional a los votos válidos & carga & ninguno \\
\addlinespace
X-01 & omitir el plano de corte cuando su resultado se descarta & votaciones & estado
ajustado sin cambio, pero altera el flujo del procedimiento: fuera del alcance sin decisión
del propietario \\
\addlinespace
P-01 & eliminar la inserción concurrente en un contenedor compartido de la fase de
legisladores & legisladores & requisito previo de corrección; no altera los números
ajustados \\
P-02 & registrar el número de hilos en el resumen de cada corrida & salida & requisito
previo de medición; la línea nueva se excluye de la comparación de archivos \\
\end{longtable}
\normalsize

\footnotesize
\begin{longtable}{L{2.2cm}L{6.2cm}L{4.0cm}}
\caption[Protocolo de medición de rendimiento.]{Protocolo de medición de la mejora de
rendimiento: qué fija cada paso y por qué.}\label{tab:protocolo}\\
\toprule
Paso & Qué se fija o se hace & Por qué \\
\midrule
\endfirsthead
\multicolumn{3}{l}{\footnotesize\emph{Cuadro \ref{tab:protocolo}, continuación}}\\
\toprule
Paso & Qué se fija o se hace & Por qué \\
\midrule
\endhead
\bottomrule
\endfoot
Requisitos previos & corregir la inserción concurrente (P-01) y registrar el número de hilos
(P-02) & un resultado no reproducible no puede medirse \\
\addlinespace
Línea base & etiquetar la versión de partida; mismo compilador, opciones, bibliotecas y máquina
para ambas versiones; registrar procesador, memoria, sistema operativo y opciones & los tiempos
publicados no fijan estas condiciones \\
\addlinespace
Paneles & cuatro estáticos y tres dinámicos del cuadro \ref{tab:paneles}; mismos valores
iniciales; cuatro ciclos; dos dimensiones & el tiempo depende del camino de búsqueda \\
\addlinespace
Condición de aceptación & cinco corridas por versión e hilos: parámetros, claves y estadísticas
idénticos; excluir solo tiempo, hilos y fecha; validar esos metadatos aparte & un cambio que
altera un bit es un cambio numérico (D-20) \\
\addlinespace
Perfilado & atribuir el tiempo de la línea base a cada rutina; contadores de acceso a memoria &
elegir cambios por evidencia y contrastar la hipótesis del recorrido \\
\addlinespace
Medición & una corrida de calentamiento descartada; al menos cinco corridas por panel, versión
y número de hilos, intercaladas; mediana, mínimo, máximo y rango intercuartílico; cambio como
cociente de medianas & separar la mejora de la variación de la máquina; comparación principal
a un hilo (RNF04) \\
\addlinespace
Orden y registro & cambios de a uno, primero los de mayor efecto esperado en la fase de
votaciones; resultados en un directorio fechado e inmutable con sus programas generadores & que
un fallo sea atribuible; trazabilidad numérica (RNF06) \\
\end{longtable}
\normalsize

\chapter{Procedencia y evidencia de esta revisión}
\label{anexo:procedencia}

Se verificaron localmente los motores en \archivo{ad945a91} y la investigación en
\archivo{3cb203e3}. Los números de línea del portal R corresponden a su archivo local modificado, SHA-256\par
{\footnotesize\texttt{170d8a4638b312e6a5d83f45d1f0ea417b843d5d50d5a6ab0fd516a767dc0090}}\par
no al blob publicado de la revisión base. La fuente original seleccionada es \archivo{DW-NOMINATE-2004.FOR};
su SHA-256 es\par
{\footnotesize\texttt{276e3cb545f2ba18d4ce4ca193f45ae8695c7e0b343f3dbcb7caaeaaa13d63a8}}\par
El ejecutable original recompilado el 6 de octubre tiene SHA-256\par
{\footnotesize\texttt{08dc4e042e46b67268a0e533ef64ba6a6643a33d6de22b6bda1d6ad13701c2ab}}\par
Los comandos, compilador GNU Fortran 15.2.0, dependencias y hashes completos están en
\href{run:../evidence/lineage/builds/README.md}{la procedencia de builds incluida}.
Esa construcción original usó los flags legacy registrados y wmay usó \texttt{-O2};
sus tiempos no se usan como benchmark. El benchmark E3 exige nuevas condiciones comparables.

Las tres comparaciones dinámicas retenidas original/wmay y sus limitaciones se documentan en
\href{run:../evidence/ENGINE-CLOSEOUT.md}{el registro de linaje incluido}. El motor fiel con
GRID activo no reemplaza esa identificación del original. El paquete conserva fuentes de
código verificadas y manifiestos; no contiene un nuevo replay numérico ni modifica el artículo.

El 8 de octubre la API confirmó la rama pública del motor en \archivo{00ba24c0}, con seis
commits locales de reorganización aún ausentes del remoto; el commit \archivo{ad945a91} no
se pudo recuperar allí. El backlog base y el portal R publicado sí existen en la revisión privada citada
de investigación, pero las propuestas I y la ampliación local del portal no están en esos blobs.
El archivo local del portal coincide con el hash usado en la revisión de figuras del 7 de octubre;
sus números de línea siguen vigentes para esa copia. Se incluye y cita esa fuente local exacta. Esto verifica destinos para la cuenta consultada, no permisos de cada lector.
La evidencia de diseño de esta revisión se entrega en este PDF y en su paquete local.

La planilla de seguimiento pudo leerse mediante el conector: G2 aún figuraba ``Sin empezar'',
con fecha de E2 del 2 de octubre; E3 y H8 no aparecían entre esas filas. El registro de este
informe documenta los artefactos actuales; sincronizar los estados y la fecha de la planilla
con el responsable sigue pendiente. El backlog incluido también conserva G2/E3 sin empezar y H8
pendiente; este informe distingue diseño escrito de aprobación. La API del 8 de octubre devuelve
\texttt{permission=none} y \texttt{pull=false} para Julio (\texttt{jcredberry}) en
\texttt{quevotan-db}; la ruta remota \texttt{tesis/} aún no existe. Publicación y acceso siguen abiertos.
La confirmación del objetivo chileno, RNF01 y la periodización siguen siendo decisiones del
profesor guía. El resultado de esta pasada es un candidato de revisión, no una entrega registrada.


\section*{Propuestas I1--I7 para la Etapa 3}

La revisión final de diseño produjo dieciocho mejoras agrupadas en siete propuestas locales
del backlog. El registro leído las marca \emph{Pendiente}, hasta que Roberto y Pablo las
adopten. No se presentan como implementación realizada ni como tareas E2 incumplidas.

\begin{table}[htbp]
\centering\footnotesize
\caption{Propuestas de la Etapa 3 incorporadas al registro local el 8 de octubre.}
\begin{tabular}{L{0.7cm}L{7.9cm}L{2.0cm}L{1.0cm}}
\toprule
ID & Alcance y relación con este diseño & Responsable & Horas \\
\midrule
I1 & corregir P-01 y verificar salidas idénticas con 1/2/4/12 hilos; prerrequisito de I6 & Pablo & 8 \\
I2 & resumen del estado final, exportación 1D y hilos en resumen & Pablo & 8 \\
I3 & errores y contratos de entrada/salida & Pablo & 8 \\
I4 & modos de prueba por componente accesibles desde la CLI fiel & Pablo & 6 \\
I5 & productor RF09 único, emparejamiento uno a uno & Roberto & 10 \\
I6 & rendimiento a un núcleo y salidas idénticas; original 2004 como referente de tesis,
wmay como antecedente adicional & ambos & 40 \\
I7 & comentarios obsoletos; rutas y dependencia requieren pruebas de ejecución y construcción & Roberto & 3 \\
\bottomrule
\end{tabular}
\end{table}

Las horas son estimaciones propuestas. La copia del backlog local incluida conserva el texto
original de I6 (``contra wmay''); el \href{run:../evidence/ETAPA-3-consistency-proposal-2026-10-08.md}{registro E3 corregido}
alinea la propuesta con el original 2004.
El registro corregido no adopta las tareas ni cambia los responsables. Las pruebas de I1 muestran
el defecto antes y exigen determinismo después; I2 corrige estadísticas antes de congelar I6.

La \href{run:../evidence/performance-candidates-confirmed-2026-10-08.md}{comprobación estática
de candidatos del 8 de octubre} confirma premisas de C1 a C15 y propone N1 a N7; no construye,
ejecuta ni mide mejoras. Precisa que C2 puede trasladar costo a legisladores y verosimilitud,
C6 consume menos resultados de lo supuesto, y C5/C8 tienen un alcance de asignaciones más
estrecho. Son insumos de I6, no resultados del informe.

Antes de optimizar, I1 debe preinicializar el mapa de periodos servidos y endurecer el acceso
al mapa de coeficientes de la misma región paralela; I2 corrige resumen, exportación y sello. Sigue una
línea base reconstruida y perfilada, preservando D-03 y verificando identidad de salidas tras
cada cambio. N1 exige reutilizar solo entradas iguales bit a bit; C14 y N4 omiten evaluaciones
del flujo heredado y requieren una decisión explícita. Los tiempos históricos citados por esa
nota deben regenerarse; no se transfieren como mediciones actuales. El catálogo del anexo
\ref{anexo:rendimiento} conserva el diseño E2 y esta fuente amplía la planificación E3.

La auditoría completa de concurrencia del motor moderno queda como tarea H4 pendiente: revisar
escrituras compartidas y repetir paneles estáticos/dinámicos a 1/2/4/12 hilos con builds e inputs
identificados. P-01 y P-02 se encontraron en el fiel; sus diferencias con el moderno no prueban
seguridad global. No se ejecutó esa auditoría ni un nuevo experimento en esta revisión.
